\section*{Bab \ref{ch:g5:natural-hazards} --- Bahaya Alam}

\begin{solution}{exo:g5:natural-hazards:1}
(a) Hanya \omterm{def:g5:natural-hazards:hazard}{bahaya alam}: tidak ada yang dicelakai. (b) \omterm{def:g5:natural-hazards:disaster}{Bencana alam}: rumah-rumah
rusak. (c) Hanya \omterm{def:g5:natural-hazards:hazard}{bahaya alam}. (d) \omterm{def:g5:natural-hazards:disaster}{Bencana alam}.
\end{solution}

\begin{solution}{exo:g5:natural-hazards:2}
Dari udara: gelombang panas, \omterm{def:g5:natural-hazards:storm}{badai}, \omterm{def:g5:natural-hazards:drought}{kekeringan}. Dari air: \omterm{def:g5:water-shapes-the-land:flood}{banjir}. Dari
tanah: \omterm{def:g5:natural-hazards:landslide}{tanah longsor}, \omterm{def:g4:volcanoes-earthquakes-first-look:eruption}{letusan} \omterm{def:g4:volcanoes-earthquakes-first-look:volcano}{gunung api}, \omterm{def:g4:volcanoes-earthquakes-first-look:earthquake}{gempa bumi}.
\end{solution}

\begin{solution}{exo:g5:natural-hazards:3}
\omterm{def:g5:natural-hazards:drought}{Kekeringan} adalah masa yang panjang, berbulan-bulan atau bertahun-tahun,
dengan \omterm{def:g4:clouds-rain-snow:precipitation}{hujan} yang jauh lebih sedikit daripada biasanya. \omterm{def:g5:natural-hazards:drought}{Kekeringan} dapat
membuat petani tidak panen, karena tanamannya layu dan mati, dan dapat
membuat kota-kota kekurangan air minum, karena \omterm{def:g1:landscapes-around-us:water}{sungai} dan \omterm{def:g1:landscapes-around-us:water}{danau} menyusut.
\end{solution}

\begin{solution}{exo:g5:natural-hazards:4}
Misalnya: memasukkan atau mengikat segala sesuatu yang dapat diterbangkan
angin (kursi taman, mainan, tempat sampah); memeriksa senter dan menyiapkan
baterai cadangan serta air; tetap di dalam rumah, jauh dari jendela, selama
\omterm{def:g5:natural-hazards:storm}{badai} berlangsung; terus mendengarkan kabar dari radio.
\end{solution}

\begin{solution}{exo:g5:natural-hazards:5}
(a) Hurikan. (b) Taifun. (c) Siklon tropis.
\end{solution}

\begin{solution}{exo:g5:natural-hazards:6}
Hijau, kuning, oranye, merah. Pada tingkat merah, bahayanya besar: segera
bertindak dan ikuti petunjuk orang yang bertanggung jawab.
\end{solution}

\begin{solution}{exo:g5:natural-hazards:7}
Sebagian lereng \omterm{def:g1:landscapes-around-us:landform}{bukit} di sebelah kanan terlepas dan meluncur ke bawah,
meninggalkan bekas luka yang gundul. \omterm{def:g3:stones-pebbles-sand:sand}{Lumpurnya} menyebar melintasi seluruh
dasar \omterm{def:g1:landscapes-around-us:landform}{lembah} dan menutupi jalan. \omterm{def:g3:stones-pebbles-sand:sand}{Lumpur} itu juga membendung \omterm{def:g1:landscapes-around-us:water}{sungai} seperti
bendungan: air \omterm{def:g1:landscapes-around-us:water}{sungai}, yang tidak dapat lagi mengalir pergi, menumpuk di
belakangnya dan menggenangi hutan di hulu.
\end{solution}

\begin{solution}{exo:g5:natural-hazards:8}
Lena yang benar. Kita tidak dapat menghentikan \omterm{def:g4:clouds-rain-snow:precipitation}{hujan} atau \omterm{def:g1:landscapes-around-us:water}{sungai}, tetapi
kita dapat menjauhkan orang dari bahaya. Misalnya: membangun rumah di tanah
yang lebih tinggi, bukan di \omterm{def:g5:water-shapes-the-land:flood}{dataran banjir}; memperingatkan orang tepat
waktu dengan peringatan \omterm{def:g5:water-shapes-the-land:flood}{banjir}, supaya mereka pindah ke tanah yang lebih
tinggi; mengajari semua orang agar tidak pernah berjalan atau berkendara
melewati air \omterm{def:g5:water-shapes-the-land:flood}{banjir}.
\end{solution}

\begin{solution}{exo:g5:natural-hazards:9}
(a) $4 \times 60 = 240$ milimeter. (b) Seperempat dari 60 adalah 15; dalam
empat bulan, $4 \times 15 = 60$ milimeter. (c) $240 - 60 = 180$ milimeter
hilang: tiga perempat dari \omterm{def:g4:clouds-rain-snow:precipitation}{hujan} yang biasa.
\end{solution}

\begin{solution}{exo:g5:natural-hazards:10}
(a) $1\,200 \div 4 = 300$ rumah. (b) Sepertiga dari 300 adalah 100; jadi
$300 - 100 = 200$ rumah masih berada di \omterm{def:g5:water-shapes-the-land:flood}{dataran banjir}.
\end{solution}

\begin{solution}{exo:g5:natural-hazards:11}
Alasan yang mungkin: desa pertama sudah diperingatkan tepat waktu dan
penduduknya sudah bersiap (daun jendela ditutup, barang-barang dimasukkan),
sedangkan desa kedua tidak; rumah-rumah di desa pertama dibangun dengan
kokoh, sedangkan rumah di desa kedua rapuh; desa kedua terletak di \omterm{def:g1:landscapes-around-us:water}{pesisir}
yang rendah atau di \omterm{def:g5:water-shapes-the-land:flood}{dataran banjir}, tempat \omterm{def:g5:natural-hazards:storm}{badai} juga membawa \omterm{def:g5:water-shapes-the-land:flood}{banjir},
sedangkan desa pertama berdiri di tanah yang lebih tinggi; di desa kedua,
pohon-pohon yang dekat dengan rumah tumbang menimpa rumah-rumah itu. Bahaya
yang sama tidak menimbulkan kerugian yang sama di mana-mana: hal itu
bergantung pada di mana dan bagaimana manusia tinggal, dan pada seberapa
baik mereka bersiap.
\end{solution}
